% ===================================================================
%  QUADRO N.º 13 — O Hotel. — O Restaurante. — O Café.
%
%  Ceci n'est PAS une transcription : c'est une traduction, faite en
%  2026, en portugais d'aujourd'hui. Voir texto/pt/10-quadro-01.tex
%  pour la consigne, qui vaut pour les seize fichiers.
%
%  « Caen » est la ville normande, et rien d'autre : les tripes se
%  font « a la moda de Caen ». Le nom reste tel quel.
% ===================================================================

%%K t13-noto-1 noto
\VUnotes{143.2mm}{%
\textit{\VUgras{(*) Para os pormenores do quarto, ver o quadro n.º 6.}}%
}

%%K t13-apar-1 apar
\VUsaut{3.0mm}
\VUtitre{15.0pt}{60}{QUADRO N.º 13}
\VUsaut{1.6mm}
\VUpk{10.2pt}{40}{O Hotel. --- O Restaurante.}
\VUsaut{2.0mm}
\VUpk{10.2pt}{40}{O Café.}
\VUsaut{3.4mm}

%%K t13-01-1 p
1. --- Tendo chegado a Royan ontem à tarde, alojei-me no \textit{Hotel do
Oceano}, que um amigo me tinha recomendado.

%%K t13-01-2 p
Deram-me um belo \VUgras{quarto} \textsuperscript{(2)} no segundo \VUgras{andar}
\textsuperscript{(1)}, onde dormi muito bem (*).

%%K t13-02-1 p
2. --- Esta manhã, ao sair do meu quarto, onde a \VUgras{criada de quarto}
\textsuperscript{(3)} me tinha trazido o pequeno-almoço, entrei na \VUgras{sala de
fumo} \textsuperscript{(5)}, que serve também de \VUgras{sala de leitura}
\textsuperscript{(4)}, para folhear os \VUgras{jornais} \textsuperscript{(6)} enquanto
fumava um \VUgras{cigarro} \textsuperscript{(8)}. Quando acabei de ler, desci no
\VUgras{elevador} \textsuperscript{(9)} e saí a dar um passeio pela cidade.

%%K t13-03-1 p
3. --- Como me tinha esquecido de perguntar ao \VUgras{rececionista}
\textsuperscript{(12)} do hotel a que horas se serviam as refeições na
\VUgras{mesa redonda} \textsuperscript{(11)}, voltei cedo e aproveitei para tomar
um banho na \VUgras{casa de banho} \textsuperscript{(10)} do hotel. Depois fui à
\VUgras{caixa} \textsuperscript{(15)} para telefonar a um dos meus amigos. Mas o
\VUgras{telefone} \textsuperscript{(13)} estava ocupado; ao ver o meu apuro, a
\VUgras{caixa} \textsuperscript{(14)}, que estava a lançar uma \VUgras{conta}
\textsuperscript{(17)} no \VUgras{registo dos viajantes} \textsuperscript{(16)}, pediu ao
\VUgras{porteiro} \textsuperscript{(18)} que mandasse o \VUgras{paquete}
\textsuperscript{(19)} fazer o meu recado na sua \VUgras{bicicleta}
\textsuperscript{(20)}.

%%K t13-04-1 p
4. --- Agradeci-lhe e saí a dar uma gorjeta ao \VUgras{carregador}
\textsuperscript{(24)} (o homem para tudo), que tinha trazido da estação no seu
\VUgras{carrinho} \textsuperscript{(25)} a minha \VUgras{chapeleira}
\textsuperscript{(26)}, a minha \VUgras{mala} \textsuperscript{(27)} e a minha
\VUgras{bolsa de viagem} \textsuperscript{(28)}. O \VUgras{carteiro}
\textsuperscript{(21)}, que chegava nesse momento, entregou-me uma \VUgras{carta}
\textsuperscript{(22)}.

%%K t13-05-1 p
5. --- Fiquei mais um pouco à soleira, junto ao \VUgras{marco do correio}
\textsuperscript{(23)}, a ver o movimento da rua. O \VUgras{condutor}
\textsuperscript{(30)} do \VUgras{ómnibus} \textsuperscript{(29)} carregava no seu
veículo a \VUgras{mala grande} \textsuperscript{(31)} de um \VUgras{viajante}
\textsuperscript{(32)} a quem o \VUgras{hoteleiro} \textsuperscript{(33)} dizia adeus. Um
jovem \VUgras{telegrafista} \textsuperscript{(35)} trazia, sem nenhuma pressa, um
\VUgras{telegrama} \textsuperscript{(34)} para um dos hóspedes do hotel; um
\VUgras{turista} \textsuperscript{(36)}, com a sua \VUgras{máquina fotográfica}
\textsuperscript{(38)} pendurada ao ombro, consultava o seu \VUgras{guia}
\textsuperscript{(37)}.

%%K t13-tit-1 sub
\VUsaut{4.0mm}
\VUpk{10.2pt}{40}{A hora do almoço.}
\VUsaut{3.0mm}

%%K t13-06-1 p
6. --- Diante da grade, um plácido \VUgras{moço de recados}
\textsuperscript{(39)} dormitava entre os seus \VUgras{ganchos de carga}
\textsuperscript{(40)} e a sua \VUgras{caixa de engraxador} \textsuperscript{(41)},
enquanto uma jovem \VUgras{lavadeira} \textsuperscript{(42)} que levava um
\VUgras{cesto} \textsuperscript{(43)} de roupa se ria do ar solene do
\VUgras{mordomo} \textsuperscript{(44)} plantado junto à porta.

%%K t13-07-1 p
7. --- A vista do \VUgras{cozinheiro} \textsuperscript{(45)}, que tinha saído da
sua \VUgras{cozinha} \textsuperscript{(47)} da \VUgras{cave} \textsuperscript{(48)} para
mandar um \VUgras{ajudante de cozinha} \textsuperscript{(46)} a um recado,
lembrou-me que se aproximava a hora do almoço, e fui ao restaurante do
outro lado do pátio, onde um \VUgras{automobilista} \textsuperscript{(50)}
estacionava o seu \VUgras{automóvel} \textsuperscript{(49)} enquanto um
\VUgras{mecânico} \textsuperscript{(52)} reparava, seguindo as indicações do
\VUgras{ciclista} \textsuperscript{(53)}, um \VUgras{triciclo a gasolina}
\textsuperscript{(51)} que acabava de ter um acidente.

%%K t13-08-1 p
8. --- Perguntei a um jovem \VUgras{groom} \textsuperscript{(54)} que levava
\VUgras{canecas de cerveja} \textsuperscript{(55)} se a mesa redonda estava sob a
varanda, e como me disse que sim, sentei-me a uma das mesas e serviram-me
uma refeição excelente.

%%K t13-noto-2 noto
\VUnotes{143.2mm}{%
\textit{\VUgras{(*) Com o auxílio da lista de pratos que o vocabulário
dá, o professor pode variar facilmente a ementa seguinte.}}%
}

%%K t13-tit-2 sub
\VUsaut{4.0mm}
\VUpk{10.2pt}{40}{Um almoço excelente.}
\VUsaut{3.0mm}

%%K t13-09-1 p
9. --- O criado trouxe-me a ementa do almoço (*) e a carta dos vinhos.
Escolhi e pedi \VUgras{camarões} com \VUgras{manteiga} de
\VUgras{entrada}, e, de primeiro prato, \VUgras{tripas à moda de Caen}.
Não tomei \VUgras{peixe}, mas pedi uma dose de \VUgras{perna} de carneiro
e, de \VUgras{legume}, \VUgras{espinafres} com natas. Depois, como nenhum
dos \VUgras{doces} me apetecia, mandei trazer um sortido de sobremesas:
\VUgras{queijo}, \VUgras{fruta} e \VUgras{bolachas}. Juntei um excelente
\VUgras{vinho} de Saint-Émilion e, muito satisfeito com a minha refeição,
levantei-me da mesa depois de dar ao \VUgras{criado} \textsuperscript{(57)} uma boa
\VUgras{gorjeta} \textsuperscript{(59)}.

%%K t13-10-1 p
10. --- Ao ver que me dispunha a partir, o \VUgras{gerente}
\textsuperscript{(60)} propôs-me sair à varanda para admirar o esplêndido
panorama do \VUgras{oceano} \textsuperscript{(62)}. Ao longe distinguiam-se
pequenos \VUgras{barcos} \textsuperscript{(63)} de pesca, \VUgras{veleiros}
\textsuperscript{(64)} e um enorme \VUgras{transatlântico} \textsuperscript{(65)} cuja
silhueta negra se recortava sobre as \VUgras{nuvens} \textsuperscript{(67)} brancas
do \VUgras{horizonte} \textsuperscript{(66)}. Uma brisa ligeira agitava a
\VUgras{bandeira} \textsuperscript{(70)} plantada por cima da \VUgras{tabuleta}
\textsuperscript{(69)} do \VUgras{vestíbulo} \textsuperscript{(68)}.

%%K t13-tit-3 sub
\VUsaut{3.0mm}
\VUpk{10.2pt}{40}{Divertimentos.}
\VUsaut{3.0mm}

%%K t13-11-1 p
11. --- O espetáculo era tão interessante que decidi subir no dia
seguinte ao terraço do café, levando comigo uns \VUgras{binóculos}
\textsuperscript{(71)} ou um \VUgras{óculo} \textsuperscript{(72)}.

%%K t13-12-1 p
12. --- Ao sair da varanda fui ao café do hotel tomar uma
\VUgras{bebida} \textsuperscript{(73)} e jogar uma partida de \VUgras{bilhar}
\textsuperscript{(75)}. Atravessei a sala do café, onde muitos clientes jogavam
\VUgras{xadrez} \textsuperscript{(78)}, \VUgras{damas} \textsuperscript{(79)},
\VUgras{dominó} \textsuperscript{(80)}, \VUgras{gamão} \textsuperscript{(81)} ou
\VUgras{cartas} \textsuperscript{(82)}, e subi direito à \VUgras{sala de bilhar}
\textsuperscript{(74)}.

%%K t13-13-1 p
13. --- Terminada a partida, desci ao rés do chão e pedi ao criado um
\VUgras{envelope} \textsuperscript{(84)}, \VUgras{papel} \textsuperscript{(83)}, um
\VUgras{selo} \textsuperscript{(85)}, uma \VUgras{caneta} \textsuperscript{(87)} e
\VUgras{tinta} \textsuperscript{(86)} para escrever uma carta.

%%K t13-13-2 p
Depois chamei um \VUgras{cocheiro} \textsuperscript{(92)} cuja \VUgras{carruagem}
\textsuperscript{(88)} tinha parado diante do café. Perguntei-lhe o seu preço
por hora e por corrida e, uma vez de acordo, abriu-me a
\VUgras{portinhola} \textsuperscript{(89)} e instalou-me lá dentro. Voltou a subir
à sua \VUgras{boleia} \textsuperscript{(91)}, tocou vivamente os cavalos com o seu
\VUgras{chicote} \textsuperscript{(93)}, e lá fomos dar um passeio delicioso.
\VUsaut{6.0mm}
\VUfilet{22mm}
